\section{Conclusões}
\label{sec:conclusoes}

Este trabalho revisou as topologias CA--CC para carregadores
de VEs e estudou o retificador Vienna a quatro fios com FP
variável. O controle proposto, com PR por fase no referencial
estacionário e referência $A\cos(\theta-\varphi)$, permite variar o FP sem
perturbar a malha de tensão. Como a ligação do neutro elimina o grau de liberdade de modo comum, a faixa de FP sem distorção fica restrita a uma janela indutiva centrada
em $\psi_0 = \arcsin(\omega_0 L A/V_{\sub{pk}})$, que se desloca com a carga, e
pedidos fora dela, principalmente capacitivos, geram distorção
e corrente de neutro elevadas. O limitador proposto mantém a
operação dentro da janela, com THD de até \SI{3,1}{\percent}, e a diferença entre o FP pedido e o obtido indica quando o pedido
excede a capacidade da topologia analisada.

